% answers_short.tex —— 2010 年·全国大纲卷Ⅰ·文科 · 简短答案速查
% 由 tools/make_short_answers.py 自 answers.tex 的 \gksolhead/\gkanswer 提取，
% 并与 2010/2010_数学真题_全国1文_答案要点.md 逐题比对后生成，勿手改（改源 MD 后重跑）。

\begin{gknotice}
\textbf{使用说明}：本节只给最终答案便于快速核对，完整推导见解析版。答案为编写者独立推导，\textbf{非官方标准答案}。
\end{gknotice}

\gkpart{选择题}
\begin{center}\begin{tabular}{ccccccccccccc}
\toprule
题号 & 1 & 2 & 3 & 4 & 5 & 6 & 7 & 8 & 9 & 10 & 11 & 12 \\
\midrule
答案 & C & C & B & A & A & C & C & B & D & C & D & B \\
\bottomrule
\end{tabular}\end{center}

\gkpart{填空题}
\begin{itemize}[leftmargin=2.6em]
  \item[\textbf{13.}] $(-1,2)\cup(2,+\infty)$
  \item[\textbf{14.}] $-\dfrac{24}{7}$
  \item[\textbf{15.}] $30$
  \item[\textbf{16.}] $\dfrac{\sqrt3}{3}$
\end{itemize}

\gkpart{解答题}
\begin{itemize}[leftmargin=2.6em]
  \item[\textbf{17.}] $S_n=\dfrac{n(3n-1)}2$ 或 $S_n=2n(5-n)$。
  \item[\textbf{18.}] $C=90^\circ$。
  \item[\textbf{19.}] (1) $\dfrac25$；(2) $\dfrac{328}{625}$
  \item[\textbf{20.}] (1) $SE=2EB$；(2) 二面角 $A-DE-C=120^\circ$
  \item[\textbf{21.}] (1) 极小值为 $-12$（在 $x=-2$ 处取得），无极大值；(2) $-\dfrac43\leq a\leq\dfrac16$
  \item[\textbf{22.}] (1) 见解析版；(2) 内切圆 $M$ 的方程为 $\left(x-\dfrac19\right)^2+y^2=\dfrac49$
\end{itemize}
